%Version 3.1 December 2024
% See section 11 of the User Manual for version history
%
%%%%%%%%%%%%%%%%%%%%%%%%%%%%%%%%%%%%%%%%%%%%%%%%%%%%%%%%%%%%%%%%%%%%%%
%%                                                                 %%
%% Please do not use \input{...} to include other tex files.       %%
%% Submit your LaTeX manuscript as one .tex document.              %%
%%                                                                 %%
%% All additional figures and files should be attached             %%
%% separately and not embedded in the \TeX\ document itself.       %%
%%                                                                 %%
%%%%%%%%%%%%%%%%%%%%%%%%%%%%%%%%%%%%%%%%%%%%%%%%%%%%%%%%%%%%%%%%%%%%%

%%\documentclass[referee,sn-basic]{sn-jnl}% referee option is meant for double line spacing

%%=======================================================%%
%% to print line numbers in the margin use lineno option %%
%%=======================================================%%

\documentclass[lineno,pdflatex,sn-basic]{sn-jnl}% Basic Springer Nature Reference Style/Chemistry Reference Style

%%=========================================================================================%%
%% the documentclass is set to pdflatex as default. You can delete it if not appropriate.  %%
%%=========================================================================================%%

%%\documentclass[sn-basic]{sn-jnl}% Basic Springer Nature Reference Style/Chemistry Reference Style

%%Note: the following reference styles support Namedate and Numbered referencing. By default the style follows the most common style. To switch between the options you can add or remove �Numbered� in the optional parenthesis. 
%%The option is available for: sn-basic.bst, sn-chicago.bst%  
 
%%\documentclass[pdflatex,sn-nature]{sn-jnl}% Style for submissions to Nature Portfolio journals
%%\documentclass[pdflatex,sn-basic]{sn-jnl}% Basic Springer Nature Reference Style/Chemistry Reference Style
%\documentclass[pdflatex,sn-mathphys-num]{sn-jnl}% Math and Physical Sciences Numbered Reference Style
%%\documentclass[pdflatex,sn-mathphys-ay]{sn-jnl}% Math and Physical Sciences Author Year Reference Style
%%\documentclass[pdflatex,sn-aps]{sn-jnl}% American Physical Society (APS) Reference Style
%%\documentclass[pdflatex,sn-vancouver-num]{sn-jnl}% Vancouver Numbered Reference Style
%%\documentclass[pdflatex,sn-vancouver-ay]{sn-jnl}% Vancouver Author Year Reference Style
%%\documentclass[pdflatex,sn-apa]{sn-jnl}% APA Reference Style
%%\documentclass[pdflatex,sn-chicago]{sn-jnl}% Chicago-based Humanities Reference Style

%%%% Standard Packages
%%<additional latex packages if required can be included here>

\usepackage{graphicx}%
\usepackage{multirow}%
\usepackage{amsmath,amssymb,amsfonts}%
\usepackage{amsthm}%
\usepackage{mathrsfs}%
\usepackage[title]{appendix}%
\usepackage{xcolor}%
\usepackage{textcomp}%
\usepackage{manyfoot}%
\usepackage{booktabs}%
\usepackage{algorithm}%
\usepackage{algorithmicx}%
\usepackage{algpseudocode}%
\usepackage{listings}%
%%%%

%%%%%=============================================================================%%%%
%%%%  Remarks: This template is provided to aid authors with the preparation
%%%%  of original research articles intended for submission to journals published 
%%%%  by Springer Nature. The guidance has been prepared in partnership with 
%%%%  production teams to conform to Springer Nature technical requirements. 
%%%%  Editorial and presentation requirements differ among journal portfolios and 
%%%%  research disciplines. You may find sections in this template are irrelevant 
%%%%  to your work and are empowered to omit any such section if allowed by the 
%%%%  journal you intend to submit to. The submission guidelines and policies 
%%%%  of the journal take precedence. A detailed User Manual is available in the 
%%%%  template package for technical guidance.
%%%%%=============================================================================%%%%

%% as per the requirement new theorem styles can be included as shown below
\theoremstyle{thmstyleone}%
\newtheorem{theorem}{Theorem}%  meant for continuous numbers
%%\newtheorem{theorem}{Theorem}[section]% meant for sectionwise numbers
%% optional argument [theorem] produces theorem numbering sequence instead of independent numbers for Proposition
\newtheorem{proposition}[theorem]{Proposition}% 
%%\newtheorem{proposition}{Proposition}% to get separate numbers for theorem and proposition etc.

\theoremstyle{thmstyletwo}%
\newtheorem{example}{Example}%
\newtheorem{remark}{Remark}%

\theoremstyle{thmstylethree}%
\newtheorem{definition}{Definition}%

\raggedbottom
%%\unnumbered% uncomment this for unnumbered level heads

\begin{document}

\title[Article Title]{Article Title}

%%=============================================================%%
%% GivenName	-> \fnm{Joergen W.}
%% Particle	-> \spfx{van der} -> surname prefix
%% FamilyName	-> \sur{Ploeg}
%% Suffix	-> \sfx{IV}
%% \author*[1,2]{\fnm{Joergen W.} \spfx{van der} \sur{Ploeg} 
%%  \sfx{IV}}\email{iauthor@gmail.com}
%%=============================================================%%

%\author*[1,2]{\fnm{First} \sur{Author}}\email{iauthor@gmail.com}
%
%\author[2,3]{\fnm{Second} \sur{Author}}\email{iiauthor@gmail.com}
%\equalcont{These authors contributed equally to this work.}
%
%\author[1,2]{\fnm{Third} \sur{Author}}\email{iiiauthor@gmail.com}
%\equalcont{These authors contributed equally to this work.}
%
%\affil*[1]{\orgdiv{Department}, \orgname{Organization}, \orgaddress{\street{Street}, \city{City}, \postcode{100190}, \state{State}, \country{Country}}}
%
%\affil[2]{\orgdiv{Department}, \orgname{Organization}, \orgaddress{\street{Street}, \city{City}, \postcode{10587}, \state{State}, \country{Country}}}
%
%\affil[3]{\orgdiv{Department}, \orgname{Organization}, \orgaddress{\street{Street}, \city{City}, \postcode{610101}, \state{State}, \country{Country}}}

%%==================================%%
%% Sample for unstructured abstract %%
%%==================================%%

\abstract{The abstract serves both as a general introduction to the topic and as a brief, non-technical summary of the main results and their implications. Authors are advised to check the author instructions for the journal they are submitting to for word limits and if structural elements like subheadings, citations, or equations are permitted.}

%%================================%%
%% Sample for structured abstract %%
%%================================%%

% \abstract{\textbf{Purpose:} The abstract serves both as a general introduction to the topic and as a brief, non-technical summary of the main results and their implications. The abstract must not include subheadings (unless expressly permitted in the journal's Instructions to Authors), equations or citations. As a guide the abstract should not exceed 200 words. Most journals do not set a hard limit however authors are advised to check the author instructions for the journal they are submitting to.
% 
% \textbf{Methods:} The abstract serves both as a general introduction to the topic and as a brief, non-technical summary of the main results and their implications. The abstract must not include subheadings (unless expressly permitted in the journal's Instructions to Authors), equations or citations. As a guide the abstract should not exceed 200 words. Most journals do not set a hard limit however authors are advised to check the author instructions for the journal they are submitting to.
% 
% \textbf{Results:} The abstract serves both as a general introduction to the topic and as a brief, non-technical summary of the main results and their implications. The abstract must not include subheadings (unless expressly permitted in the journal's Instructions to Authors), equations or citations. As a guide the abstract should not exceed 200 words. Most journals do not set a hard limit however authors are advised to check the author instructions for the journal they are submitting to.
% 
% \textbf{Conclusion:} The abstract serves both as a general introduction to the topic and as a brief, non-technical summary of the main results and their implications. The abstract must not include subheadings (unless expressly permitted in the journal's Instructions to Authors), equations or citations. As a guide the abstract should not exceed 200 words. Most journals do not set a hard limit however authors are advised to check the author instructions for the journal they are submitting to.}

\keywords{keyword1, Keyword2, Keyword3, Keyword4}

%%\pacs[JEL Classification]{D8, H51}

%%\pacs[MSC Classification]{35A01, 65L10, 65L12, 65L20, 65L70}

\maketitle

\section{Introduction}\label{sec1}

The Introduction section, of referenced text \cite{bib1} expands on the background of the work (some overlap with the Abstract is acceptable). The introduction should not include subheadings.

Springer Nature does not impose a strict layout as standard however authors are advised to check the individual requirements for the journal they are planning to submit to as there may be journal-level preferences. When preparing your text please also be aware that some stylistic choices are not supported in full text XML (publication version), including coloured font. These will not be replicated in the typeset article if it is accepted. 

\section{Results}\label{sec2}

\subsection{Twin system architecture}
% 介绍我们这个孪生系统的组成结构，系统运行流程，侧重点在于系统的整体布局，突出整体性
% 介绍系统的运行流程
This system consists of data collection nodes deployed at five major intersections (each equipped with one LiDAR, six cameras, and seven V2X devices) and a centralized server (equipped with one V2X device and the CARLA simulation platform). It operates using a pipeline mechanism that involves “real-time data collection and transmission, online inference, and streaming simulation.”
During operation, each collection endpoint first captures point clouds and images. After preprocessing using PV-RCNN, YOLOv8, and ResNet-50 to extract the targets’ 3D and 2D bounding boxes and appearance feature vectors, the data is transmitted to the server via V2X.
After the server has collected all the data, it uses the JPDA algorithm to perform multi-object tracking at a single intersection and infers trajectories in unknown areas using a fifth-degree polynomial, and then completes cross-intersection trajectory association based on feature matching.
Finally, the system smooths and denoises the associated trajectories and uses a vector velocity control strategy based on adjacent sampled frames to drive the target in segmented, constant-speed linear motion, thereby reproducing the spatio-temporally consistent twin trajectories with high precision in the CARLA virtual environment.

We conducted experiments on four different maps to test the effectiveness of this system.
Among these, Town01 and Town10 are maps included in the CARLA simulation software; both feature standard road topologies, spatial divisions for mixed-use traffic (pedestrians and vehicles), and complete 3D environments that include buildings and traffic infrastructure.
HutbCarlaCity and Zhongdian are maps we created based on real-world road networks and environments; the former depicts Hunan University of Commerce and its surrounding areas, while the latter depicts the Changsha China Electronics Software Park and its surrounding areas.
We ran 1,200 frames on each map to simulate traffic conditions over a two-minute period, and used the overall average across five intersections to represent the final result for that map.


\subsection{System physical layout and transmission latency}
% 介绍我们系统的数据是怎样端到端传输的，收集端与服务端的布局是怎样的，如何计算延迟，详细说明”在线“属性
% 本论文的特色，创新点
Our system consists of a data collection end and a central server.
We selected five major intersections on each map as data collection points.
One lidar sensor, six cameras, and seven V2X communication devices have been deployed at each intersection.
We deployed the LiDAR at the center of the intersection, 1.8 meters above the ground.
The six cameras, centered around the LiDAR unit, are positioned as follows: front, rear, left front, left rear, right front, and right rear. They are all uniformly mounted at a height of 3.6 meters above the ground, spaced 14 meters apart front to back and 8 meters apart left to right.
The seven V2X communication devices are each mounted 0.8 meters above the lidar and the six cameras.
The central server is located at the point closest to all intersections to minimize data transmission latency.
At the same time, deploy one V2X communication device at the central server to receive data transmitted from the data collection endpoints.

The latency in this system primarily stems from three sources: processing latency at the data collection end, data transmission latency, and processing latency at the central server.
Processing latency at the data collection end primarily stems from the preprocessing of raw data. We mark the moment just before the data enters preprocessing as the start time and the moment preprocessing ends as the end time, then calculate the difference between the two as the processing latency.
To calculate data transmission latency, we use the V2X transmission time at the data collection end as the start time and the V2X reception time at the central server as the end time, and calculate the difference between the two as the transmission latency.
Since the central server must receive data from all intersections before it can perform trajectory tracking, the longest transmission delay is therefore the total transmission delay.
Processing latency on the central server primarily stems from the execution time of the relevant programs, namely the generation and re-creation of trajectories.
We use the time when the data is input as the start time and the time when the target is successfully replicated in the virtual world as the end time, and calculate the difference between the two as the processing delay.
Total delay refers to the time elapsed from when the target appears in front of the sensor until it is successfully detected.


\begin{table}[htbp]
  \centering
  \caption{System latency.}
  \label{tab:algorithm_performance}

  \begin{tabular*}{\textwidth}{@{\extracolsep{\fill}}ccccc}
    \toprule
    Scene & Data Preprocessing & Data Transmission & Trajectory Tracking and Replay & Total \\
    \midrule
      Town01  & 2 & 1 & 1 &                    \\
      Town10  & 400 & 533 & 600 & 1533                   \\
    \bottomrule
  \end{tabular*}
\end{table}


\subsection{Generation of twin trajectories}
% 介绍我们系统是如何生成孪生轨迹的，轨迹的生成方法
% 分单路口轨迹与多路口轨迹，未知区域轨迹预测
% 说明用到的公式，模型，以及其在孪生系统中的作用，突出系统的创新性

\begin{table}[htbp]
  \centering
  \footnotesize
  \setlength{\tabcolsep}{3pt}
  \caption{System latedncy.}
  \label{tab:tracking-results}

  \begin{tabular*}{\textwidth}{@{\extracolsep{\fill}}cccccccccccc}
    \toprule
    \textbf{Scene} & \textbf{Cat.} & \textbf{Rcll} & \textbf{Prcn} & \textbf{FTR} & \textbf{FP} & \textbf{FN} & \textbf{IDS} & \textbf{MT} & \textbf{ML} & \textbf{MOTA} & \textbf{MOTP} \\
    \midrule
    \multirow{2}{*}{Town01} & person & 77.55 & 97.95 & 0.07 & 35.8 & 411.8 & 14.8 & 72.95 & 16.19 & 75.05 & 92.47 \\
                            % & person & 77.5512 & 97.9486 & 0.0716 & 35.8 & 411.8 & 14.8 & 72.9524 & 16.1905 & 75.0471 & 92.4724 \\
                            & vehicle & 83.48 & 84.89 & 0.54 & 267.6 & 332.2 & 10.8 & 81.31 & 13.69 & 68.30 & 87.75 \\
    \midrule
    \multirow{2}{*}{Town10} & person & 80.50 & 76.27 & 2.71 & 1353.6 & 1011 & 70.6 & 61.31 & 10.58 & 53.14 & 91.12 \\
                            & vehicle & 88.17 & 71.19 & 2.30 & 1153.6 & 320 & 22.2 & 83 & 11.43 & 45.47 & 85.64 \\
    \midrule
    \multirow{2}{*}{HutbCarlaCity} & person &  &  &  &  &  &  &  &  &  &  \\
                            & vehicle &  &  &  &  &  &  &  &  &  &  \\
    \midrule
    \multirow{2}{*}{Zhongdian} & person &  &  &  &  &  &  &  &  &  &  \\
                            & vehicle &  &  &  &  &  &  &  &  &  &  \\
    \bottomrule
  \end{tabular*}
\end{table}

To demonstrate the effectiveness of this system, we conducted validation tests across four different scenarios.
Town01 and Town10 are virtual scene maps.
HutbCarlaCity and Zhongdian are real-world maps.
We ran 500 frames for each scenario to simulate traffic conditions over a period of time, and used the overall average across the five intersections to represent the final results, yielding the experimental results shown in Table \ref{tab:tracking-results}.

In a 500-frame simulation, this system successfully achieved a spatio-temporally consistent digital twin.
From Table \ref{tab:tracking-results}, we can see that:
In the virtual scenarios, the MOTA (75.05 and 68.30) and Prcn (97.95 and 84.89) for the Town01 scenario were generally higher than those for the Town10 scenario (MOTA of 53.14 and 45.47, Prcn of 76.27 and 71.19) . The highest MOTA and Prcn values both occurred in the pedestrian class in the Town01 scenario (75.05 and 97.95, respectively), while the lowest MOTA value occurred in the vehicle class in the Town10 scenario (45.47).
In both scenarios, the recall rates for vehicles (83.48 for Town01 and 88.17 for Town10) and MT (81.31 for Town01 and 83 for Town10) were both higher than those for pedestrians in the same scenarios. The highest recall rate was observed in the vehicle class for Town10 (88.17).
All error metrics (FP, FN, IDS, FTR) for Town10 are significantly higher than those for Town01. Specifically, the “human” class in Town10 has the highest values for FP (1,353.6), FN (1,011), IDS (70.6), and FTR (2.71); the “human” class in Town01 has the lowest FP value (35.8).


\subsection{Twins in different scenarios}
% 对比分析虚拟场景与真实场景下的孪生效果，说明系统的鲁棒性，突出系统的适应性
% 通过表格的形式列出实验结果进行对比分析，得出结论
% 孪生效果包括跟踪精度，延迟，最终复现轨迹的精度


\subsection{System robustness testing and sensor network optimization}
% 测试各种极端条件(如大雾、暴雨等)对孪生的影响，找出对孪生效果起决定性作用(基本可以认为无法孪生，即孪生精度过低)的那个临界值；考虑传感器数目、位置等对孪生效果的影响分布
% 通过表格的形式列出实验结果进行对比分析，得出结论
% 孪生效果包括跟踪精度，延迟，最终复现轨迹的精度


\section{Discussion}\label{sec12}
% 结论部分介绍孪生的作用，意义，启发
% 不足及未来可能的改进方向
% 介绍本论文的创新点

Discussions should be brief and focused. In some disciplines use of Discussion or `Conclusion' is interchangeable. It is not mandatory to use both. Some journals prefer a section `Results and Discussion' followed by a section `Conclusion'. Please refer to Journal-level guidance for any specific requirements. 



%%===========================================================================================%%
%% If you are submitting to one of the Nature Portfolio journals, using the eJP submission   %%
%% system, please include the references within the manuscript file itself. You may do this  %%
%% by copying the reference list from your .bbl file, paste it into the main manuscript .tex %%
%% file, and delete the associated \verb+\bibliography+ commands.                            %%
%%===========================================================================================%%

\bibliography{sn-bibliography}% common bib file
%% if required, the content of .bbl file can be included here once bbl is generated
%%\input sn-article.bbl

\end{document}
